% Generated draft: requires mathematical and source review.
\documentclass[11pt]{article}
\usepackage[margin=25mm]{geometry}
\usepackage{fontspec}
\setmainfont{Latin Modern Roman}
\usepackage{amsmath,amssymb,mathtools,unicode-math}
\setmathfont{Latin Modern Math}
\usepackage{enumitem,needspace,xurl,hyperref}
\setlength{\parindent}{0pt}
\setlength{\parskip}{5pt}
\setlength{\emergencystretch}{3em}
\newcommand{\N}{\mathbb{N}}
\newcommand{\Z}{\mathbb{Z}}
\newcommand{\Q}{\mathbb{Q}}
\newcommand{\R}{\mathbb{R}}
\newcommand{\C}{\mathbb{C}}
\newcommand{\F}{\mathbb{F}}
\newcommand{\A}{\mathbb{A}}
\newcommand{\PP}{\mathbb{P}}
\newcommand{\E}{\mathbb{E}}
\newcommand{\eps}{\varepsilon}
\newcommand{\dd}{\,\mathrm{d}}
\begin{document}
\section*{UnsolvedMath ID: 20000990}
\textbf{Source problem code:} catalogue entry
\subsection{Definitions and mathematical statement}
Andrea Sportiello asks for a higher-dimensional explorer model and whether it is more amenable to generalization than the abelian sandpile model. The source gives no axioms for ``higher-dimensional'' or criterion for comparing amenability. The aid identifies reference [16] as the Diaconis--Fulton explorer model, now called multiple-source internal aggregation or IDLA.

A precise stack-aggregation framework for the aid's proposed construction is this. Let $\Sigma$ be a finite state set, $R\subseteq\Sigma$ absorbing roots, and let each $s\notin R$ have an infinite stack of destinations along a specified transition graph. Start with finitely many particles. Each nonroot state has capacity one; when it contains at least two particles, a legal firing sends one particle to the next destination in its stack. Particles arriving at roots are absorbed. A stable configuration has at most one particle at each nonroot state. Call stacks fair when every allowed outgoing destination occurs infinitely often in its source stack. Assume every nonroot state reaches a root in the allowed transition graph.

The aid claims that this construction terminates for fair stacks, that the stable configuration and stack-use odometer do not depend on legal firing order, and that initial particle count equals final nonroot particle count plus root absorptions. It proposes applying this framework to a finite pure simplicial complex $X$ of dimension $d\geq1$, with states the oriented $(d-1)$-faces and transitions from a Parzanchevski--Rosenthal $p$-lazy oriented ridge walk. The precise kernel, orientation convention, admissible $p$, and roots must be specified before this is a fully determined model; they are not given in the record. The aid reports specialization to the graph explorer model in dimension one, subject to those state and occupancy conventions.

For the version sharing a single occupancy slot between opposite orientations, the aid reports an exact factorization through ordinary IDLA on the unsigned upper-coface adjacency graph: its vertices are unoriented $(d-1)$-faces, with adjacency when two distinct such faces belong to a common $d$-simplex, with the projected walk's transition probabilities. This quotient may lose topology. The filled triangle, viewed as a pure $2$-complex, has three edge states and adjacency graph $C_3$; its boundary cycle, viewed as a pure $1$-complex, has three vertex states and the same graph. With corresponding lazy transition and root conventions, the resulting explorer processes agree, while
\[
 H_1(\text{filled triangle};\Z)=0,\qquad
 H_1(\text{boundary cycle};\Z)\cong\Z.
\]
This supplied obstruction shows that this particular quotient process need not detect homology. It does not rule out other higher-dimensional explorer models or settle the workshop's comparative question.
\subsection{Short English statement}
Seek a higher-dimensional extension of explorer aggregation. The supplied aid proposes abelian stack aggregation on oriented faces, but its one-slot unoriented quotient becomes ordinary graph IDLA and can erase homology. Model conventions and the intended standard of higher-dimensional behavior remain unresolved.
\subsection{Sources}
\begin{itemize}
\item AIM workshop \emph{Generalizations of chip-firing and the critical group} (2013), problem 25, Andrea Sportiello: \url{https://aimath.org/pastworkshops/chipfiringproblems.pdf}.
\item The original cites reference [16]; the attached aid identifies it as Diaconis--Fulton explorer aggregation and names the Parzanchevski--Rosenthal oriented codimension-one-face walk. Full bibliographic entries and the walk kernel are absent.
\item Local registry \texttt{data/unsolvedmath/problems.json}, permanent ID 20000990, exact source code \texttt{catalogue entry}; canonical location \texttt{aim-combinatorics-notes.json notes[114]}. Attached attempt 1, partial result, problem status at run open, rigor proved, low novelty confidence; machine-generated, not peer reviewed, requiring independent verification before citation.
\end{itemize}
\end{document}
